\subsection{Duality of finitely generated modules}

In this section we consider a Noetherian ring A of finite dimension which possesses a dualizing module \(\Omega\) . One then has \(\mathrm{di}_{\Lambda}(\Omega) = \dim(\mathrm{A}) < +\infty\) ( \(n^{\circ}\) 1, prop. 1). Choose an injective resolution of finite length \(e: \Omega \to (\mathbf{I}, \delta)\) . For every complex C of A-modules, denote \(\mathbf{D}(\mathbf{C})\) the complex \(\operatorname{Homgr}_{\Lambda}(\mathbf{C}, \mathbf{I})\) . This applies in particular to every A-module M, considered as a complex concentrated in degree 0; we then have \(\mathbf{D}(\mathrm{M})^{i} = \mathrm{Hom}_{\mathrm{A}}(\mathrm{M},\mathrm{I}^{i})\) for every integer \(i\) . Recall that in A, X, p.~100, th. 1, we constructed a canonical isomorphism

\[
\varphi (\mathrm{M}, \mathrm{I}): \mathrm{H} (\mathbf {D} (\mathrm{M})) \longrightarrow \operatorname{Ext} _ {\mathrm{A}} (\mathrm{M}, \Omega).
\]

Examples.--- 1) The complex \(\mathbf{D}(\mathbf{A}) = \operatorname{Homgr}_{\mathbf{A}}(\mathbf{A},\mathbf{l})\) is identified with I. The map \(c:\Omega \to \mathbf{D}(\Lambda)\) is by definition a homologism.

\begin{enumerate}
\def\labelenumi{\arabic{enumi})}
\setcounter{enumi}{1}
\item
  The homomorphism \(e \in \text{Homgr}_A(\Omega, I)^0\) is an element of \(\mathbf{D}(\Omega)^0\) ; the A-linear map \(\tilde{e}: A \to \mathbf{D}(\Omega)\) such that \(\tilde{e}(1) = e\) is a homologism (no. 4, prop. 7, a) and b)).
\item
  Let S be a multiplicative subset of A. The \(S^{-1}A\) -module \(S^{-1}\Omega\) is dualizing (no. 1, cor. 1 of prop. 2); the \(S^{-1}A\) -modules \(S^{-1}I^{i}\) are injective (cor. 1 of prop. 3 of § 3, no. 2) and the morphism \(e^{\prime}:S^{-1}\Omega\to S^{-1}I\) deduced from e is an injective resolution of \(S^{-1}\Omega\) , to which one may therefore apply the preceding results. For every complex C of finite type (and in particular every finitely generated A-module M), the canonical homomorphism from \(S^{-1}\mathbf{D}(\mathbf{C})\) in \(\operatorname{Homgr}_{\mathrm{S}^{-1}\mathrm{A}}(\mathrm{S}^{-1}\mathrm{C},\mathrm{S}^{-1}\mathrm{I})=\mathbf{D}(\mathrm{S}^{-1}\mathrm{C})\) is bijective.
\item
  Let A be a Dedekind ring, K its field of fractions. The A-module A is dualizing and admits the injective resolution I of length 1 defined by the exact sequence
\end{enumerate}

\[
0 \to \mathrm{A} \stackrel {e} {\longrightarrow} \mathrm{K} \stackrel {\delta} {\longrightarrow} \mathrm{K/A} \to 0
\]

where δ is the canonical surjection. For every A-module M, the complex D(M) is the complex concentrated in degrees 0 and 1

\[
\dots \longrightarrow 0 \longrightarrow \operatorname{Hom} _ {\mathrm{A}} (\mathrm{M}, \mathrm{K}) \stackrel {d} {\longrightarrow} \operatorname{Hom} _ {\mathrm{A}} (\mathrm{M}, \mathrm{K} / \mathrm{A}) \longrightarrow 0 \longrightarrow \dots
\]

with \(d = \mathrm{Hom}_{\mathrm{A}}(1_{\mathrm{M}},\delta)\) . We have an exact sequence

\[
0 \to \operatorname{Hom} _ {\mathrm{A}} (\mathrm{M}, \mathrm{A}) \longrightarrow \mathbf {D} (\mathrm{M}) ^ {0} \stackrel {d} {\longrightarrow} \mathbf {D} (\mathrm{M}) ^ {1} \longrightarrow \operatorname{Ext} _ {\mathrm{A}} ^ {1} (\mathrm{M}, \mathrm{A}) \to 0.
\]

For every morphism of complexes \(f: \mathrm{C} \to \mathrm{C}'\) , we denote \(\mathbf{D}(f): \mathbf{D}(\mathrm{C}') \to \mathbf{D}(\mathrm{C})\) the morphism of complexes \(\operatorname{Homgr}_{\mathrm{A}}(f, 1_{\mathrm{I}})\) If \(f\) is a homologism, \(\mathbf{D}(f)\) is a homologism (A, X, p.~86, prop. 4, b)). If \(\mathrm{C}' \xrightarrow{f} \mathrm{C} \xrightarrow{g} \mathrm{C}''\) is an exact sequence of complexes, the sequence of complexes \(\mathbf{D}(\mathrm{C}''') \xrightarrow{\mathbf{D}(g)} \mathbf{D}(\mathrm{C}) \xrightarrow{\mathbf{D}(f)} \mathbf{D}(\mathrm{C}')\) is exact (A, X, p.~83, prop. 2, a)).

Let M be an A-module. To each element m of M, associate the map \(\alpha_{\mathrm{M}}(m): f \mapsto f(m)\) of D(M) into I; it is an element of \(\mathbf{D}(\mathbf{D}(\mathbf{M}))_{0} = \operatorname{Homgr}_{\mathrm{A}}(\mathbf{D}(\mathbf{M}), \mathrm{I})_{0}\) It follows from the definitions that \(\alpha_{\mathrm{M}}(m)\) is a morphism of complexes, hence an element of \(\mathrm{Z}_{0}(\mathbf{D}(\mathbf{D}(\mathbf{M})))\) .

One thus defines a morphism of complexes:

\[
\alpha_ {\mathrm{M}}: \mathrm{M} \rightarrow \mathbf {D} (\mathbf {D} (\mathrm{M})) ,
\]

whence, by passing to homology, a homomorphism of A-modules

\[
\alpha_ {\mathrm{M}}: \mathrm{M} \to \mathrm{H} _ {0} (\mathbf {D} (\mathbf {D} (\mathrm{M}))).
\]

THEOREM 1. Let M be a finitely generated A-module. Then \(\alpha_{\mathrm{M}}\) is a homomorphism: we have \(\mathrm{H}_{i}(\mathbf{D}(\mathbf{D}(\mathbf{M}))) = 0\) for \(i \neq 0\) and the homomorphism \(\alpha_{\mathrm{M}}\) is bijective.

Let us first take M = A. The map \(e : \Omega \to \mathbf{D}(A)\) is a homomorphism (example 1), hence so is the map \(\mathbf{D}(e) : \mathbf{D}(\mathbf{D}(A)) \to \mathbf{D}(\Omega)\) . The map \(\tilde{e} : A \to \mathbf{D}(\Omega)\) is a homologism (example 2), and we have \(\mathbf{D}(e) \circ \alpha_{\mathbf{A}} = \tilde{c}\) ; thus \(\alpha_{A}\) is a homomorphism, which proves the theorem in this case. It follows that \(\alpha_{M}\) is a homomorphism when the A-module M is finitely generated free.

Let us turn to the general case; we shall prove by induction on the integer n the following assertion:

\((A_{n})\) for every finitely generated A-module M, the homomorphism \(H_{i}(\alpha_{M})\) is bijective for \(i \leqslant n\) .

This also means that \(\mathrm{H}_{i}(\mathbf{D}(\mathbf{D}(\mathbf{M})))\) is zero for \(i \neq 0\) and \(i \leqslant n\) , and that \(\alpha_{M}\) is bijective if \(n \geqslant 0\) . Observe that \((\mathrm{A}_{n})\) holds for n \textless{} -d, where d is the length of the complex 1: indeed the A-module \(\mathbf{D}(\mathbf{D}(\mathbf{M}))_{i}\) is equal to \(\bigoplus_{p}\mathrm{Hom}_{\mathrm{A}}(\mathrm{Hom}_{\mathrm{A}}(\mathrm{M},\mathrm{I}^{p}),\mathrm{I}^{p-i})\) , hence is zero for i \textless{} -d and i \textgreater{} d.

Let us prove the implication \((\mathrm{A}_{n}) \Rightarrow (\mathrm{A}_{n+1})\) . Let M be a finitely generated A-module. There exists a finitely generated free A-module L and an exact sequence \(0 \to N \xrightarrow{u} L \xrightarrow{v} M \to 0\) . The sequence \(0 \to \mathbf{D}(M) \xrightarrow{\mathbf{D}(v)} \mathbf{D}(L) \xrightarrow{\mathbf{D}(u)} \mathbf{D}(N) \to 0\) is exact; likewise, if we set \(u' = \mathbf{D}(\mathbf{D}(u))\) and \(v' = \mathbf{D}(\mathbf{D}(v))\) , the sequence \(0 \to \mathbf{D}(\mathbf{D}(N)) \xrightarrow{u'} \mathbf{D}(\mathbf{D}(L)) \xrightarrow{v'} \mathbf{D}(\mathbf{D}(M)) \to 0\) is exact.

Since \(H_{i}(\mathbf{D}(\mathbf{D}(\mathrm{L})))\) is zero for \(i \neq 0\) we have isomorphisms

\[
\mathrm{H} _ {i} (\mathbf {D} (\mathbf {D} (\mathrm{M}))) \longrightarrow \mathrm{H} _ {i - 1} (\mathbf {D} (\mathbf {D} (\mathrm{N}))) \quad \text { pour } i \neq 0, 1;
\]

this entails the implication \((\mathrm{A}_{n})\Rightarrow(\mathrm{A}_{n+1})\) for \(n\neq-1\) and \(n\neq0\) . Consider the commutative diagram with exact rows

\[
\begin{array}{c c c c c c c} 0 \to & N & \xrightarrow {u} & L & \xrightarrow {v} & M & \to 0 \\ & \alpha_ {N} \Bigg \downarrow & & \alpha_ {L} \Bigg \downarrow & & \alpha_ {M} \Bigg \downarrow \\ 0 \to H _ {1} (\mathbf {D} (\mathbf {D} (M))) & \xrightarrow {} H _ {0} (\mathbf {D} (\mathbf {D} (N))) & \xrightarrow {H _ {0} (u ^ {\prime})} & H _ {0} (\mathbf {D} (\mathbf {D} (L))) & \xrightarrow {H _ {0} (v ^ {\prime})} & H _ {0} (\mathbf {D} (\mathbf {D} (M))) \longrightarrow H _ {- 1} (\mathbf {D} (\mathbf {D} (N)) \end{array}
\]

where \(\alpha_{\mathrm{L}}\) is bijective. If \((\mathrm{A}_0)\) is satisfied, the homomorphism \(\alpha_{\mathrm{N}}\) is also bijective, hence \(\mathrm{H}_0(u')\) is injective and one obtains \(\mathrm{H}_1(\mathbf{D}(\mathbf{D}(\mathbf{M}))) = 0\) , whence \((\mathrm{A}_1)\) If \((\mathrm{A}_{-1})\)

is satisfied, \(\mathrm{H}_{-1}(\mathbf{D}(\mathbf{D}(\mathrm{N})))\) is zero, therefore \(\mathrm{H}_{0}(v^{\prime})\) is surjective, which implies that \(\alpha_{M}\) is surjective. Since this holds for every finitely generated A-module M, \(\alpha_{N}\) is also surjective; by I, § 1, no. 4, cor. 2 of prop. 2, \(\alpha_{M}\) is bijective, so that \((\mathrm{A}_{0})\) is satisfied.

Thus \((\mathrm{A}_{n})\) is true for all n, which proves the theorem.

Let M be a finitely generated A-module; set \(c = \dim(A) - \dim_{A}(M)\) . Let \(\mathbf{D}'(M)\) the subcomplex of \(\mathbf{D}(M)\) equal to \(\bigoplus_{i\in\mathcal{O}}\mathbf{D}(M)^{i}\bigoplus Z^{c}(\mathbf{D}(M))\) , and by

\[
j: \mathbf {D} ^ {\prime} (\mathrm{M}) \rightarrow \mathbf {D} (\mathrm{M})
\]

the canonical injection. From the canonical surjection we deduce \(\mathrm{Z}^{c}(\mathbf{D}(\mathrm{M}))\to\mathrm{H}^{c}(\mathbf{D}(\mathrm{M}))\) and of the isomorphism \(\varphi(\mathrm{M},\mathrm{I})\) a morphism of complexes

\[
p: \mathbf {D} ^ {\prime} (\mathrm{M}) (- c) \rightarrow \operatorname{Ext} _ {\Lambda} ^ {c} (\mathrm{M}, \Omega);
\]

as \(\mathrm{H}^{i}(\mathbf{D}(\mathrm{M}))\) is zero for i \textless{} c (no. 1, prop. 3 a)), \((\mathbf{D}^{\prime}(\mathrm{M})(-c), p)\) is a left resolution of \(\operatorname{Ext}_{\mathrm{A}}^{c}(\mathrm{M}, \Omega)\) . By A, X, p.~100, th. 1, we have a canonical isomorphism

\[
\varphi^ {0} \left(\mathbf {D} ^ {\prime} (\mathrm{M}) (- c), \mathrm{I}\right): \mathrm{H} _ {0} \left(\mathbf {D} \left(\mathbf {D} ^ {\prime} (\mathrm{M})\right)\right) \longrightarrow \operatorname{Ext} _ {\mathrm{A}} ^ {c} \left(\operatorname{Ext} _ {\mathrm{A}} ^ {c} (\mathrm{M}, \Omega), \Omega\right);
\]

composing it with the homomorphism \(\mathrm{H}_{0}(\mathbf{D}(j)) : \mathrm{H}_{0}(\mathbf{D}(\mathbf{D}(\mathrm{M}))) \to \mathrm{H}_{0}(\mathbf{D}(\mathbf{D}'(\mathrm{M})))\) we therefore obtain a homomorphism

\[
\mathrm{H} _ {0} (\mathbf {D} (\mathbf {D} (\mathrm{M}))) \longrightarrow \operatorname{Ext} _ {\mathrm{A}} ^ {c} \left(\operatorname{Ext} _ {\mathrm{A}} ^ {c} (\mathrm{M}, \Omega), \Omega\right),
\]

whence finally, by composition with \(\alpha_{\mathrm{M}}\) , a canonical homomorphism

\[
\beta_ {\mathrm{M}}: \mathrm{M} \longrightarrow \operatorname{Ext} _ {\mathrm{A}} ^ {c} (\operatorname{Ext} _ {\mathrm{A}} ^ {c} (\mathrm{M}, \Omega), \Omega).
\]

COROLLARY.--- If the A-module M is Macaulay, the homomorphism \(\beta_{\mathrm{M}}\) is bijective.

If M is Macaulay, the \(\Lambda\) -module \(\mathrm{H}^i(\mathbf{D}(\mathrm{M}))\) is zero for \(i \neq c\) ( \(n^\circ 1\) (cor. of prop. 3), so that the canonical injection \(j: \mathbf{D}'(\mathrm{M}) \to \mathbf{D}(\mathrm{M})\) is a homologism; consequently the morphism of complexes \(\mathbf{D}(j): \mathbf{D}(\mathbf{D}(\mathrm{M})) \to \mathbf{D}(\mathbf{D}'(\mathrm{M}))\) is a homologism ( \(\Lambda\) , X, p.~86, prop. 4). Thus \(\mathrm{H}_0(\mathbf{D}(j))\) is bijective; on the other hand, \(\alpha_{\mathrm{M}}\) is bijective by th. 1, whence the corollary.

When the A-module M is of finite length, the A-module \(\operatorname{Ext}_{\Lambda}^{c}(M,\Omega)\) is identified with the Matlis dual of M (cf.~§ 8, n° 5, example 3 and th. 3), and one recovers prop. 4 of § 8, n° 4.

